% Created 2026-04-15 mié 21:28
% Intended LaTeX compiler: pdflatex
\documentclass[11pt]{article}
\usepackage[utf8]{inputenc}
\usepackage[T1]{fontenc}
\usepackage{graphicx}
\usepackage{longtable}
\usepackage{wrapfig}
\usepackage{rotating}
\usepackage[normalem]{ulem}
\usepackage{amsmath}
\usepackage{amssymb}
\usepackage{capt-of}
\usepackage{hyperref}
\author{Abraham Nuñez}
\date{\today}
\title{Notas Aprendizajemaquina}
\hypersetup{
 pdfauthor={Abraham Nuñez},
 pdftitle={Notas Aprendizajemaquina},
 pdfkeywords={},
 pdfsubject={},
 pdfcreator={Emacs 30.2 (Org mode 9.7.39)}, 
 pdflang={English}}
\begin{document}

\maketitle
\tableofcontents

\section{Informacion}
\label{sec:org4f7593c}
\begin{itemize}
\item Maestro: David Martinez Galicia
\item Materia: Aprendizaje Maquina
\item Carrera: CIberseguridad e Infraestructura de Computo
\item Contacto: davidmartinez02@uv.mx
\end{itemize}
\section{Informacion}
\label{sec:orgb0a7e25}
\subsection{Evaluacion}
\label{sec:orgf0ff750}
\begin{itemize}
\item Examenes: 40\%
\item Tareas: 30\%
\item Proyecto: 30\%
\end{itemize}
\section{Temario}
\label{sec:orge341fb2}
\begin{itemize}
\item 1. Introduccion
\item 2. Manejo de datos
\item 3. Aprendizaje supervisado
\item 4. Aprendizaje no supervisado
\item 5. Aprendizaje por refuerzo
\item 6. Evaluacion de modelos
\end{itemize}
\section{Introduccion}
\label{sec:org384c1c5}
\subsection{Inteligencia Artificial}
\label{sec:orgb6d3a3e}
\begin{itemize}
\item Que es la inteligencia artificial?
\item ``Cada aspecto de la inteligencia artificial puede ser descrito con suficiente precision que puede fabricarse una maquina para simularlo''.
\item El concepto existe desde 1956.
\item Inteligencia artificial: Un enfoque moderno, Pearson Rusell, S. J., y Norving P. (2021).
\item Concuerdan que la IA debe estudiar sistemas inteligentes.
\item Caracteristicas para ser inteligentes:
\end{itemize}
\subsubsection{Cognicion y comportamiento}
\label{sec:orgf729da0}
\begin{itemize}
\item Pensar (crear) -> Humano
\item Actuar (simular) -> Racional
\end{itemize}
\subsection{Tipos de sistemas}
\label{sec:orgd41c00c}
\begin{itemize}
\item Sistemas que piensan como humanos
\item Sistemas que piensan racionalmente
\item Sistemas que actuan como humanos
\item Sistemas que actuan racionalmente
\end{itemize}
\subsection{Simulacion del comportamiento humano}
\label{sec:orge064234}
\begin{itemize}
\item La prueba de Turing
\item El sistema debe de poseer las siguientes capacidades:
\begin{itemize}
\item Procesamiento de lenguaje natural
\item Representacion del conocimiento
\item Razonamiento automatico
\item Aprendizaje maquina
\end{itemize}
\item Para superar la prueba de Turing la computadora debe incluir:
\begin{itemize}
\item Vision computacional para percibir objetos.
\item Robotica para manipular y mover objetos.
\end{itemize}
\end{itemize}
\subsection{Aprendizaje Maquina}
\label{sec:orgd5f2089}
\begin{itemize}
\item Se refiere a un espectro de situaciones en las cuales un agente incrementa su conocimiento para cumplir una tarea.
\item Machine Learning. Mitchell, T. M. (1997).
\item Un programa de computadora aprende de la experiencia E, con respecto a alguna clase de tareas T y una medida de desempeno P, si su desempeno en las tareas T, medido por P, mejora con la experiencia E.
\end{itemize}
\subsubsection{Tipos de aprendizajes}
\label{sec:org89c0f34}
\begin{itemize}
\item Aprendizaje supervisado: Aprender una funcion a partir de ejemplos de sus entradas y sus salidas.
\begin{itemize}
\item Clasificacion: Diagnostico de enfermedades.
\item Regresion: Prediccion de la temperatura.
\end{itemize}
\item Aprendizaje no supervisado: Aprender a partir de patrones de entradas que no tienen valores de salidas.
\begin{itemize}
\item Agrupacion: Segmentacion de clientes.
\item Reduccion de dimensionalidad: Visualizar datos.
\item Estimacion de la densidad: Describir la altura de una poblacion.
\end{itemize}
\item Aprendizaje por refuerzo: Aprender una forma de actuar usando un refuerzo (recompens o penalizacion)
\begin{itemize}
\item Aprender politicas de accion: Actuar en un entorno.
\end{itemize}
\end{itemize}
\subsection{Introduccion al manejo de datos}
\label{sec:org631d826}
\begin{itemize}
\item La calidad de los datos determina la calidad del modelo.
\item ``Garbage in, garbage out''.
\item Sin la preparacion adecuada de los datos, no hay aprendizaje efectivo.
\item Transformamos datos crudos en patrones para la toma de decisiones.
\item Datos limpios dectetan amenazas reales en ciberseguridad.
\end{itemize}
\subsubsection{Datos estructurados}
\label{sec:org5ed7743}
\begin{itemize}
\item Datos organizados en formato tabular con esquema definido.
\end{itemize}
\begin{enumerate}
\item Caracteristicas:
\label{sec:orgb09a81d}
\begin{itemize}
\item Formate predecible.
\item Tipos de datos definidos.
\item Faciles de almacenar y consultar en bases de datos relacionales.
\item Ejemplos:
\begin{itemize}
\item Tablas SQL.
\item Hojas de calculo Excel/CSV.
\item Logs de sistemas con formatos fijo.
\end{itemize}
\end{itemize}
\end{enumerate}
\subsubsection{Datos NO estructurados}
\label{sec:org454e43f}
\begin{itemize}
\item Datos sin formato predefinido ni organizacion tabular.
\end{itemize}
\begin{enumerate}
\item Caracteristicas:
\label{sec:orgf41adc4}
\begin{itemize}
\item Sin esquema fijo.
\item Requieren procedimiento previo para extraer informacion.
\item Representan 80\% de los datos empresariales,
\item Ejemplos:
\begin{itemize}
\item Texto.
\item Imagenes y videos.
\item Audio y grabaciones.
\end{itemize}
\end{itemize}
\end{enumerate}
\subsubsection{Datos etiquetados}
\label{sec:org62e848c}
\begin{itemize}
\item Conjunto de datos donde cada ejemplo tiene una etiqueta/clase asociada.
\end{itemize}
\begin{enumerate}
\item Caracteristicas:
\label{sec:orgd0149be}
\begin{itemize}
\item Requieren proceso manual o semi-automatico de etiquetado.
\item Esenciales para entrenar modelos supervisados.
\item Costosos de obtener y validar.
\item Ejemplos:
\begin{itemize}
\item Correos spam/no spam
\item Imagenes de rostros con nombre de persona
\end{itemize}
\end{itemize}
\end{enumerate}
\subsubsection{Datos NO etiquetados}
\label{sec:orgfcffd6f}
\begin{itemize}
\item Datos sin categorias o respuestas conocidas.
\end{itemize}
\begin{enumerate}
\item Caracteristicas:
\label{sec:org9862a7b}
\begin{itemize}
\item Abundantes y faciles de recolectar.
\item Requieren tecnicas de descubrimiento de patrones.
\item Utiles para deteccion de anomalias.
\item Ejemplos:
\begin{itemize}
\item Trafico de red
\item Comportamiento de usuario
\item Codigo malicioso
\end{itemize}
\end{itemize}
\end{enumerate}
\subsubsection{Repositorio de datos}
\label{sec:org3b87195}
\subsubsection{Formatos de datos de texto}
\label{sec:org9de4795}
\begin{itemize}
\item Los formatos de texto permiten almacenar y procesar datos eficientemente.
\item Son los mas comunes para cargar en la mayoria de las herramientas
\end{itemize}
\begin{enumerate}
\item Tipos de datos de texto
\label{sec:orgf4a1ac3}
\begin{itemize}
\item CVS: Formato tabular simple y universalmente compatible.
\item TSV: Similar al CSV, pero con tabulaciones como delimitador.
\item JSON: Ligero yflexible para datos estructurados y anidados.
\item XML: Formato estructurado y autodescriptivo, aunque verboso.
\item ARFF: Incluye metadatos, usado comunmente en Weka.
\item Texto plano: Sin formato, para datos sin estructura o documentos.
\end{itemize}
\end{enumerate}
\subsubsection{Limpieza y procesamiento de datos}
\label{sec:org510d824}
\begin{itemize}
\item Los datos del mundo real son inexactos, inconsistentes e incompletos.
\item El procesamiento de datos asegura la integridad de los datos.
\item Permite generar analisis confiables y conclusiones validas.
\item Evita decisiones incorrectas o modelos ineficaces.
\end{itemize}
\begin{enumerate}
\item FLujo de trabajo de limpieza y preprocesamiento
\label{sec:orgd50a2f4}
\begin{itemize}
\item Recopilacion: Obtener datos de bases de datos, APIs o scraping.
\item Limpieza: Corregir errores, eliminar duplicados y manejar valores faltantes.
\item Integracion: Unir datos de multiples fuentes, resolviendo inconsistencias.
\item Transformacion: Codificar, normalizar y crear nuevas variables.
\item Reduccion: Seleccionar caracteristicaspara simplificar el analisis.
\end{itemize}
\end{enumerate}
\subsubsection{Python para el aprendizaje maquina}
\label{sec:org35bfaa0}
\begin{itemize}
\item Lenguaje preferido por su facilidad y bibliotecas especializados.
\end{itemize}
\begin{enumerate}
\item Bibliotecas importantes
\label{sec:orge80feb7}
\begin{itemize}
\item Pandas: Manipulacion de datos.
\item Numpy: Calculos numericos.
\item SciPy: Computo cientifico, estadistica y senales.
\item Matplotlib/Seaborn: Visualizacion de datos.
\item Scikit-learn: Machine learning y preprocesamiento.
\end{itemize}
\end{enumerate}
\subsection{Técnicas de codificación}
\label{sec:orgf88b41d}
\begin{itemize}
\item Las variables categóricas representan grupos categóricas.
\item Algunos modelos de aprendizaje maquina requieren entradas numéricas.
\item Transformar categorías en números para análisis y modelado.
\end{itemize}
\subsubsection{Tecnicas comunes:}
\label{sec:org28b193d}
\begin{itemize}
\item One-Hot Encoding: Para categorías sin orden (ej: puerto de embarque).
\item Label Encoding: Para categorías con orden (ej: nivel educativo).
\item Puntuación Z: Es una medida estadistica que indica cuantas desviaciones estandar se encuentra por encima o por debajo de la medida de una distribución.
\end{itemize}
\subsection{Aprendizaje supervisado}
\label{sec:orgd554a07}
\begin{itemize}
\item Los humanos aprendemos de experiencias pasadas.
\item Un nino aprende a identificar animales viendo ejemplos.
\item La experiencia se acumula y permite generalizar.
\item Una computadora no tiene expericencias. En su lugar, aprende de datos.
\item Estos datos son ejemplos historicos con entradas y salidas conocidas.
\item El objetivo es obtener una funcion que prediga salidas para nuevas entradas.
\end{itemize}
\subsubsection{Flujo de aprendizaje supervisado}
\label{sec:org1bb89bb}
\begin{itemize}
\item Para generar un modelo se necesita un conjunto de datos, un algoritmo de aprendizaje y una forma de evaluarlo o interpretarlo.
\begin{itemize}
\item Conjunto de datos
\item Modelo supervisado: Existen mas de cientos cincuenta modelos clasificadores:
\begin{itemize}
\item Regresion lineal.
\item Redes neuronales.
\item Redes bayesianas.
\item Arboles de decision.
\end{itemize}
\item Interpretacion/Evaluacion: Una vez obtenido el modelo, se puede:
\begin{itemize}
\item Analizar el conocimiento o los patrones capturados.
\item Medir su desempeno en datos no observados antes.
\end{itemize}
\end{itemize}
\end{itemize}
\begin{enumerate}
\item Datos en el aprendizaje supervisado
\label{sec:orgd57ac84}
\begin{itemize}
\item Un conjunto de datos compuesto por:
\begin{itemize}
\item Una lista de registros (observaciones o ejemplos)
\item Un conjunto de variables predictorias (caracteristicas o atributos)
\begin{itemize}
\item Las variables cuantitativas se conocen como continuas o numericas.
\item Las variables cualitativas se conocen como discretas o categoricas.
\item Dentro de las categorias (cualitativas) existen las nominales y las ordinales.
\end{itemize}
\item Una variable por predecir (dependiente, objetivo o clase).
\end{itemize}
\item Si la variable objetivo es categorica se vuelve un problema de clasificacion.
\item Si la variable objetivo es numerica se vuelve un problema de regresion.
\end{itemize}
\item Modelos supervisados de aprendizaje
\label{sec:orgaea1cb2}
\begin{itemize}
\item Un modelo es una funcion o teoria que se aprenda de los datos.
\item Puede verse como una hipotesis sobre la relacion entre entradas y salidas.
\item Se ajuta (entrena) un modelo usando un algoritmo de aprendizaje.
\item El algorimo (metodo o tecnica) encuentra los \textbf{parametros} del modelo.
\item Los algoritmos se configuran a traves de sus \textbf{hiperparametros}.
\item Una vez entrenado, el modelo puede predecir salidas para nuevos datos.

\item Si la variable objetivo es categorica, el modelo se llama \textbf{clasificador}.
\item Si la variable objetivo es numerica, el modelo se llama de \textbf{regresion}.
\end{itemize}
\item Interpretacion y evaluacion de modelos.
\label{sec:org1fb6cb6}
\begin{itemize}
\item Los modelos no solo predicen, tambien representan conocimiento.
\item Ese conocimiento revela patrones sobre como se relacionan las variables.
\item Es importante para tomar decisiones en situaciones complejas.
\item Algunos modelos permiten interpretar ese conocimiento.
\item Los modelos que no se pueden interpretar se denominan de caja negra.
\item Evaluar un modelo es tan importante como entrenarlo.
\item La evaluacion determina si el modelo es confiable o debe sustituirse.
\end{itemize}
\item Modelo de regresion lineal
\label{sec:org25d0f9e}
\begin{itemize}
\item Describe la relacion entre la variable independiente y una dependiente.
\item Asume que dicha relacion puede representarse mediante una linea recta.
\item Permite predecir valores continuos a partir de un solo atributo.
\item Es el puntaje de partida para entender las tareas de regresion.
\item y = mx + b (Funcion de una recta).
\item y = O(0) + 0(1)\textsuperscript{x}
\begin{itemize}
\item y: valor predicho por el modelo
\item x: variable independiente
\item O(0): intercepto que representa el valor de y cuando x = 0
\item O(1): pendiente que representa el cambio en y por cada unidad en x.
\end{itemize}
\item El objetivo del modelo es que y sea lo mas cercano posible a y.
\item El modelo define una familia de todas las lineas posibles.
\item Los parametros 0 = [O(0), O(1)]\textsuperscript{T} determinen la linea especifica.
\item Funcion de perdida (Loss Function): Mide que tan mal se ajusta el modelo a los datos.
\begin{itemize}
\item Para cada punto, calculamos el error: diferencia entre el valor real y predicho.
\item Se usa el error cuadratico para que la direccion no importe.
\item La perdida total es la suma de errores cuadraticos (SSE).
\end{itemize}
\end{itemize}
Agregar una imagen con la formula
\end{enumerate}
\end{document}
